\section{Implementation Details} \label{apdx:impl-details}
We conduct the cold-start training with the LlamaFactory \footnote{https://github.com/hiyouga/LlamaFactory} and all the RL experiments based on the verl \footnote{https://github.com/volcengine/verl} framework. 

\textbf{Cold-start. }
For the cold-start stage, we collect 40k training data from the responses of Doubao-Seed-1.6-thinking \footnote{https://openrouter.ai/bytedance-seed/seed-1.6}, including 20k math reasoning trajectories and 20k rubricator trajectories, which are collected via reject sampling. 
We then perform full-parameter supervised fine-tuning (SFT) on the 40k dataset using Qwen3 prompt template, with a sequence cutoff length of 32768, per-device train batch size = 1, gradient accumulation steps = 1, base learning rate = 2e-5 (cosine scheduler with min lr = 2e-6), weight decay = 0.1, warmup ratio = 0.01 and total training epochs = 5.0. 
For the verifier, we distill a small-sized model Qwen3-4B-Base based on the verification responses from Doubao-Seed-1.6-thinking, using the same SFT recipe as in the cold-start stage.


\begin{table*}[h]
\centering
\begin{tabularx}{0.55\linewidth}{l|c}
  \toprule
\multicolumn{1}{c}{\textbf{Hyperparameter}} \vline & \multicolumn{1}{c}{\textbf{Value}} \\
\midrule
\small
Prompt template                       & Qwen3 \\
Sequence cutoff length                & 32768 \\
Per-device train batch size           & 1 \\
Gradient accumulation steps           & 1 \\
Learning rate (cosine scheduler)      & 2e-5 (min lr = 2e-6) \\
Weight decay                          & 0.1 \\
Warmup ratio                          & 0.01 \\
Epochs                                & 5.0 \\
\bottomrule
\end{tabularx}
\caption{Key training hyperparameters used in the cold-start SFT.}
\label{tab:sft_hyper}
\end{table*}

\clearpage
\textbf{RL training. }
For the RL training, we train all the models with the DAPO-Math-17K data, with prompt max length = 16384, response max length = 12288 (overlong buffer enabled, length = 4096), no KL reward/loss, clip ratio = 0.2, learning rate for actor = 1e-6, learning rate for critic = 1e-5, temperature = 1.0, top\_p = 1.0, train batch size = 32, rollout number $N$ = 8, mini batch size = 64, weight decay=0.1, grad clip=1.0, lr warmup steps=10. 
For all RL training, we fix the max training steps to 1500.
The response length and prompt length are configured in this setting because the outputs generated by the reasoner will serve as the input prompts for the rubricator. Therefore, the max prompt length should be greater than the response length, and the sum of these two values must not exceed 32k tokens. 


\begin{table*}[h]
\centering
\begin{tabularx}{0.4\linewidth}{l|c}
  \toprule
\multicolumn{1}{c}{\textbf{Hyperparameter}} \vline & \multicolumn{1}{c}{\textbf{Value}} \\
\midrule
\small
Algorithm                             & PPO \\
Training dataset                      & DAPO-Math-17K \\
Max prompt length                     & 16384 \\
Max response length                   & 12288 \\
Overlong buffer                       & 4096 \\
KL reward/loss                        & None \\
Clip ratio                            & 0.2 \\
Learning rate (actor)                 & 1e-6 \\
Learning rate (critic)                & 1e-5 \\
Sampling temperature                  & 1.0 \\
Top\_p                                & 1.0 \\
Train batch size                      & 32 \\
Rollout number ($N$)                  & 8 \\
Mini batch size                       & 64 \\
Weight decay                          & 0.1 \\
Gradient clip                         & 1.0 \\
LR warmup steps                       & 10 \\
Max steps                             & 1500 \\
\bottomrule
\end{tabularx}
\caption{Key training hyperparameters used in RL training.}
\label{tab:rl_hyper}
\end{table*}

\clearpage
\section{Prompts} \label{apdx:prompts}
In this section, we report the prompts used as the reasoner $\pi_{\theta}^{\text{\textcolor{reasoner}{$Rea$}}}$ and the rubricator $\pi_{\theta}^{\text{\textcolor{rubricator}{$Rub$}}}$. It is worth mentioning that during our initial attempts to directly prompt Doubao-Seed-1.6-thinking model to generate rubrics based on the question and a reference answer, the model was unable to generate rubrics that are not already presented in the reference. Therefore, we modified the prompt to instruct the model to generate a set of rubrics where the current response score falls below the middle value, and explicitly reinforced this in the format reward. We observed that this practice triggers the model to actively refine rubrics during reasoning to generate meaningful rubrics for room of self-improvement. The 20k rubricator trajectories were then collected using this specific prompt with rejection sampling under such format reward.



\begin{figure}[!ht]
    \centering
    \begin{lrbox}{\DSRPrompt}
    \begin{gptpromptbox}
    \begin{minipage}{0.95\linewidth}
    \raggedright
    % \small\ttfamily
    \scriptsize\ttfamily
\setlength{\baselineskip}{1\baselineskip}
\{question\}

Please reason step by step, and put your final answer within \textbackslash
 boxed\{\}.
\end{minipage}
    \end{gptpromptbox}
    \end{lrbox}
    \usebox{\DSRPrompt}
    \caption{The prompt of the reasoner role.}
    \label{fig:evaluation-prompts}
\end{figure}


\begin{figure}[ht]
    \centering
    \begin{lrbox}{\DSRPrompt}
    \begin{gptpromptbox}
    \begin{minipage}{0.95\linewidth}
    \raggedright
    % \small\ttfamily
    \scriptsize\ttfamily
\setlength{\baselineskip}{1\baselineskip}
You are an expert in educational assessment and rubric design.
Your task is to analyze a given question–answer pair and generate comprehensive evaluation rubrics that can be used to assess response quality for this question. The answer is only a reference answer from the student(not necessarily a good response), so your rubric system should consider the merits already presented in the response, and most importantly, room for further improvements.

\# Input Data

[Question]:
\{question\}

[Response]:
\{response\}


\# Task Instructions

Based on the provided question and answer, generate a comprehensive rubric that's suitable with multiple evaluation criteria. Each criterion should be:\\
1. Specific and Measurable: Clearly define what constitutes meeting or not meeting the criterion\\
2. Binary Evaluable: Can be assessed as true/false by an LLM evaluator\\
3. Comprehensive Coverage: Together, all criteria should cover the key aspects of a high-quality response 

\# Required Rubric Categories

Generate criteria to cover these aspects:
1. Effectiveness of Problem Decomposition \& Planning
Definition: The systematic breakdown of a complex problem into logically ordered, manageable sub-tasks, and the formulation of a strategic roadmap for solving them.
Core Aspects: Identifying key components, establishing step dependencies, sequencing operations, and anticipating challenges before execution.

2. Effectiveness of BackTracking / Self-Validation / Error Handling
Definition: The continuous monitoring of reasoning validity through consistency checks, proactive error detection, and adaptive revision of flawed steps.
Core Aspects: Implementing verification checkpoints (e.g., unit analysis/estimation), diagnosing inconsistencies, recovering from dead ends via path correction (not restarting), and mitigating error propagation.

3. Reasoning Clarity \& Flow
Definition: The coherent structuring and communication of logical progression where each step explicitly follows from prior conclusions and leads to subsequent insights.
Core Aspects: Hierarchical argument organization under clear objective, unambiguous terminology, explicit logical connectors (e.g., "thus," "since," "conversely"), and seamless transitions between ideas.

4. Reasoning Focus \& Efficiency
Definition: The maintenance of strict alignment with problem objectives while optimizing cognitive resources through elimination of redundancies and irrelevant explorations.
Core Aspects: Sustained direction toward core goals, pruning of tangential paths, and avoidance of overthinking/repetition during reasoning.

5. Other Question-Specific Aspects
Definition: The strategic selection and application of domain-relevant methods that align with both problem constraints and solution goals.
Core Aspects: Leveraging context-optimal approaches (e.g., dimensional analysis in physics, elimination methods for multiple-choice), avoiding technique mismatches, and adapting tools to exploit problem structure (e.g., symmetry in geometry).


\# Output Format
Return a JSON object with the following structure:

```json\\
\{\\
  \quad"question\_domain": "math/calculus",\\
  \quad "rubrics": [
    \{
      "category": "Effectiveness of Problem Decomposition \& Planning" 
      "criterion": "XXX",
      "points": 5
    \},
    \{
      "category": "Reasoning Clarity \& Flow" 
      "criterion": "XXX",
      "points": 3
    \},
    \{
      "category": "Reasoning Focus \& Efficiency" 
      "criterion": "XXX",
      "points": 4
    \},
    \{
      "category": "Other Question-Specific Aspects" 
      "criterion": "XXX",
      "points": -3
    \},
    \{
      "category": "Reasoning Focus \& Efficiency" 
      "criterion": "XXX",
      "points": -2
    \}
    // ... additional criteria
  ],\\
  \quad "maximum\_score": a(sum of all positive metrics),\\
  \quad "minimum\_score": b(sum of all negative metrics),\\
  \quad "current\_score": x(must be smaller than (a+b)/2, otherwise revise your rubric system to get better room of improvement) \\
\}\\
```

\# Important Guidelines 
- First output the domain of the question in two levels, then generate 5 - 15 criteria in total, ensuring comprehensive coverage - Points should reflect the relative importance of each criterion (supports positive scores from 1 to 5 for reward criteria, and negative scores from -5 to -1 for penalty criteria). Specifically, a penalty criteria should be a statement that expresses flaw/failure in the reasoning process, while a reward criteria should be a statement that express merits/goal accomplishment in the reasoning process. Some of the final things that your should CHECK BEFORE YOUR FINAL RESPONSE:

\quad - Try to make your rubrics specific to the question, but not so detailed as some of the overly detailed metrics may not apply to all the response(e.g. "fail to apply XXX theorem" would be a acceptable metric, but "fail to reflect after miscalculating using XXX theorem in step25" would be a very bad metric to avoid, since another students response may not apply XXX theorem, nor having step 25). In other words, do not let the reference answer constrain your rubric system!  Again, DO NOT let the reference answer constrain your rubric system!!! To do this, try to propose a set of metrics pretending the reference answer is not there, and reconsider adding some of them. 

\quad - Remember, the answer above is only a reference answer from the student(not necessarily a good response), so your rubric system should consider the merits already presented in the response, and most importantly, room for further improvements. To enforce this, the reference response should have a score strictly below the average of your proposed rubric system. In order to check for this, you need to calculate the maximum\&minimum score from you rubric system, and score the given response at the end of your json. If you found that the rubrics you presented the first time gives a overly-high score, brainstorm more rubrics pretending the reference answer is not there in orderly to find room for improvements and add them to your rubric systems.
- Return only the JSON object without additional commentary.
    
    \end{minipage}
    \end{gptpromptbox}
    \end{lrbox}
    \usebox{\DSRPrompt}
    \caption{The prompt of the rubricator role.}
    \label{fig:evaluation-prompts}
\end{figure}


\begin{figure}[ht]
    \centering
    \begin{lrbox}{\DSRPrompt}
    \begin{gptpromptbox}
    \begin{minipage}{0.95\linewidth}
    \raggedright
    % \small\ttfamily
    \scriptsize\ttfamily
\setlength{\baselineskip}{1\baselineskip}
You are an expert in evaluating student responses to math problems. Your task is to assess a given response based on a set of binary rubrics and compute a final score.

[Task Input]:

- Question: The math problem the student is solving.

- Response: The student's step-by-step solution.

- Rubrics: A list of criteria for evaluation. Each criterion has:

\quad - category: The aspect of reasoning being assessed.
    
\quad - criterion: A binary statement (either a merit to reward or a flaw to penalize).
    
\quad - points: The points assigned (positive for rewards, negative for penalties).

[Task Instruction]:\\
Evaluate each criterion:\\
- For each criterion in the rubrics list (in the given order), determine if the criterion is satisfied by the response.\\
\quad - Output True if the criterion is met (for a reward) or if the flaw is present (for a penalty).\\
\quad - Output False otherwise. \\
- Base your judgement solely on the content of the response and the specific wording of the criterion.\\
- Compute the final score: Sum the points of every criterion that is evaluated as True (this includes both positive and negative points).

[Output Format]\\
Return a JSON object with the following structure:\\
```json\\
\{\\
  \quad "judgement": [boolean, boolean, ...],  // List of True/False for each criterion in the same order as the rubrics\\
  \quad "final\_score": number                  // Sum of points for True judgements\\
\}\\
```

[Important Guidelines]

- Be objective: Do not consider any external knowledge beyond the provided question, response, and rubrics.

- Binary decision: Each criterion must be evaluated as strictly True or False. There are no partial credits.

- Order matters: The judgement list must exactly match the order of the rubrics provided.

- Independent evaluation: Assess each criterion separately. The presence or absence of one flaw/merit does not influence the evaluation of others (unless the criterion explicitly references another).

\quad- Penalty criteria: For a criterion with negative points, True means the flaw is present (so the negative points are added). False means the flaw is absent (so 0 is added for that criterion).

\quad- Reward criteria: For a criterion with positive points, True means the merit is present (so positive points are added). False means the merit is absent (so 0 is added).

Now, proceed to evaluate the following:

Question:
\{question\}

Response:
\{resp\}

Rubrics:
\{rubrics\}

Your final output should contain only the JSON object. Do not include any additional text or explanations.

\end{minipage}
    \end{gptpromptbox}
    \end{lrbox}
    \usebox{\DSRPrompt}
    \caption{The prompt of the verifier.}
    \label{fig:evaluation-prompts}
\end{figure}

\clearpage
\section{Case Study} \label{apdx:case-study}
In this section, we present a case to illustrate how the rubricator proposes meaningful metrics conditioned on the reasoner's response. 
This example shows that the rubricator successfully identifies the key problem(\textit{Employs ad-hoc manual listing for small digit sets instead of systematic permutation formula application, increasing error risk}) that caused miscalculation for $d=0$ and $d=1$ subsets, leading to the final error of 131. It also gives a hint to efficiently solving the problem and reducing reasoning redundancy (\textit{Avoids redundant recounting by categorizing sets with minimum permutation $>$ 2013 upfront, reducing per-set analysis effort}). Therefore, the self-proposed rubrics can effectively guide the model to solve the specific problem, serving as valuable auxiliary reward for the reasoner.

\textbf{Question:}
2013 is the first year since the Middle Ages that consists of four consecutive digits. How many such years are there still to come after 2013 and before 10000?

\textbf{Ground Truth:}
149

\textbf{Prediction:}
131



\textbf{Generated Rubrics:}
\begin{figure}[!ht]
    \centering
    \begin{lrbox}{\DSRPrompt}
    \begin{promptbox}
    \raggedright
    % \small\ttfamily
    \scriptsize\ttfamily
\setlength{\baselineskip}{1\baselineskip}
\{\\
\hspace{2em}"question\_domain":"math/digit-sequence-analysis",\\
\hspace{2em}"rubrics": [\\
\hspace{4em}\{\\
\hspace{6em}"category":"Effectiveness of Problem Decomposition \& Planning",\\
\hspace{6em}"criterion":"Explicitly outlines digit set ranges, leading zero constraints, and year boundary limits before permutation calculations",\\
\hspace{6em}"points":3\\
\hspace{4em}\},\\
\hspace{4em}\{\\
\hspace{6em}"category":"Reasoning Clarity \& Flow",\\
\hspace{6em}"criterion":"Uses hierarchical structuring (e.g., bullet points) for each digit set with explicit permutation counts and exclusion justifications",\\
\hspace{6em}"points":3\\
\hspace{4em}\},\\
\hspace{4em}\{\\
\hspace{6em}"category":"Reasoning Focus \& Efficiency",\\
\hspace{6em}"criterion":"Avoids redundant recounting by categorizing sets with minimum permutation > 2013 upfront, reducing per-set analysis effort",\\
\hspace{6em}"points":3\\
\hspace{4em}\},\\
\hspace{4em}\{\\
\hspace{6em}"category":"Effectiveness of BackTracking, Self-Validation \& Error Handling",\\
\hspace{6em}"criterion":"Fails to implement edge case validation (e.g., overcounting leading-zero permutations) via secondary checklist reviews",\\
\hspace{6em}"points":-4\\
\hspace{4em}\},\\
\hspace{4em}\{\\
\hspace{6em}"category":"Other Question-Specific Aspects",\\
\hspace{6em}"criterion":"Employs ad-hoc manual listing for small digit sets instead of systematic permutation formula application, increasing error risk",\\
\hspace{6em}"points":-2\\
\hspace{4em}\}\\
\hspace{2em}],\\
\hspace{2em}"maximum\_score":9,\\
\hspace{2em}"minimum\_score":-6\\
\}
    \end{promptbox}
    \end{lrbox}
    \usebox{\DSRPrompt}
    \caption{The rubricator's generated rubrics}
    \label{fig:evaluation-prompts}
\end{figure}


\clearpage
\textbf{Reasoner's Response:}

\begin{figure}[!ht]
    \centering
    \begin{lrbox}{\DSRPrompt}
    \begin{promptbox}
    \raggedright
    % \small\ttfamily
    \scriptsize\ttfamily
\setlength{\baselineskip}{1\baselineskip}


Got it, let's try to figure out this problem: we need to find how many years after 2013 and before 10000 have four consecutive digits. First, let's make sure I understand what "four consecutive digits" means in a year (which is a four-digit number, right? Wait, years like 2013 are four-digit, 10000 is five-digit). So a four-digit year with four consecutive digits—wait, does that mean the digits are consecutive in order, like 2013 is 2,0,1,3? Wait no, 2013 isn't consecutive digits. The problem says "2013 is the first year since the Middle Ages that consists of four consecutive digits"—wait, maybe it's digits that are consecutive integers, but not necessarily in order? Wait, wait, let's recheck: "consists of four consecutive digits"—probably, it means the set of digits is four consecutive integers, maybe in any order? Wait, 2013: digits are 2,0,1,3—are those four consecutive? 0,1,2,3 are four consecutive, but 2013 has 0, so 0 is there, then 1,2,3 but missing 2? Wait no, 0,1,2,3 are four consecutive numbers, but 2013 has digits 2,0,1,3—wait maybe I misread: 2013's digits: 2,0,1,3—wait 0,1,2,3 are four consecutive, but is 2013 a permutation of four consecutive digits? Let's list permuted four-digit numbers with digits being four consecutive integers.\\
First, what are four consecutive integers? Let's denote four consecutive digits as d, d+1, d+2, d+3 where d is from 0 to 6 (since d+3 $\leq$ 9, because digits go up to 9). So possible d: 0 (0-3), 1 (1-4), 2 (2-5), 3 (3-6), 4 (4-7), 5 (5-8), 6 (6-9). Wait let's list the sets:\\
d=0: \{0,1,2,3\}\\
d=1: \{1,2,3,4\}\\
d=2: \{2,3,4,5\}\\
d=3: \{3,4,5,6\}\\
d=4: \{4,5,6,7\}\\
d=5: \{5,6,7,8\}\\
d=6: \{6,7,8,9\}\\
These are the possible sets of four consecutive digits (digits are 0-9, so four consecutive can't include 7-10, etc.). Now, a four-digit year can't start with 0, so we need to count permutations of each set that are four-digit numbers (so first digit $\ne$ 0) and are between 2013 and 10000 (so starting from 2014 up to 9999). Wait 2013 is the first since Middle Ages—so we need years after 2013, so starting from 2014.\\
Let's take each set and count valid permutations (four-digit, no leading zero, >2013):\\
\hl{Set \{0,1,2,3\} (d=0): How many four-digit permutations? Total permutations: 4! = 24, minus those starting with 0: 3! = 6, so 18. Now we need permutations >2013. Let's list or calculate: First digit: can be 1,2,3 (since 0 is invalid). If first digit is 1: then the rest are permutations of 0,2,3: 3! = 6 numbers (1023, 1032, 1203, 1230, 1302, 1320). The smallest in this first digit group is 1023, which is less than 2013, so we can ignore all first digit 1 permutations. If first digit is 2: then permutations starting with 2: the remaining digits are 0,1,3. We need permutations >2013. The smallest permutation starting with 2 is 2013, which is the given year. So permutations starting with 2 and >2013: remaining digits 0,1,3—how many are >2013? The smallest is 2013 (which is equal), next would be 2031, 2103, 2130, 2301, 2310. Wait let's count: For first digit 2, second digit: can be 0,1,3 (since 2 is already used). If second digit is 0: then third and fourth digits are permutations of 1,3: 2 numbers (0132 invalid? Wait no, first digit 2, second 0: 2 0 x x—so numbers are 2013, 2031 (since 1 and 3 are digits, so 2013 and 2031). 2013 is equal, so 2031 is >2013. If second digit is 1: then third and fourth digits are permutations of 0,3: numbers are 2103, 2130 (both >2013). If second digit is 3: then third and fourth digits are permutations of 0,1: numbers are 2301, 2310 (both >2013). So for first digit 2: number of permutations >2013: let's count: 2031 (yes), 2103 (yes), 2130 (yes), 2301 (yes), 2310 (yes). That's 5 numbers from set \{0,1,2,3\}.
Set \{1,2,3,4\} (d=1): Four-digit numbers, no leading zero (so first digit 1-4). All permutations are valid since first digit can't be 0. Total permutations: 4! = 24. We need years after 2013. The smallest number in this set is 1234, which is less than 2013, so many numbers here are below 2013. Wait first digit 1: numbers start with 1, so 1xxx, which are all <2000, so ignore. First digit 2: permutations starting with 2: total 6 (since second digit: 3! for remaining 3 digits). These are from 2134 upwards? Wait list them: 2134, 2143, 2314, 2341, 2413, 2431. All these are >2013, right? Yes, since smallest is 2134, which is >2013. So all 6 permutations starting with 2 in this set are valid.}\textbf{==> Employs ad-hoc manual listing for small digit sets instead of systematic permutation formula application, increasing error risk}\\
\hl{Set \{2,3,4,5\} (d=2): First digit can be 2-5 (since no 0 here). So permutations: total 4! =24, minus leading zero? Wait no, no 0, so all 24 are valid four-digit (since first digit 2-5). We need years after 2013. Let's list first digits: 2,3,4,5. First digit 2: permutations starting with 2: 6 numbers (2 \_ \_ \_). The smallest is 2345? Wait no, digits are 2,3,4,5—wait permutations starting with 2: second digit can be 3,4,5 (since 2 is used). So numbers: 2345, 2354, 2435, 2453, 2534, 2543. The smallest in first digit 2 is 2345, which is greater than 2013. So all 6 permutations starting with 2 are valid. First digit 3: permutations starting with 3: 6 numbers (all >2345, so >2013), same for 4 and 5: all 6 each. So total for this set: 6 + 6 + 6 + 6 = 24.\\
Set \{3,4,5,6\} (d=3): Similar to above, digits 3-6, no 0/1. So all four-digit permutations are valid (first digit 3-6). All numbers here are $\geq$ 3456, which is way above 2013, so all 4! =24 permutations are valid.\\
Set \{4,5,6,7\} (d=4): Even higher, 4xxx and above, all permutations valid: 24 numbers.\\
Set \{5,6,7,8\} (d=5): 5xxx and above: 24 numbers.\\
Set \{6,7,8,9\} (d=6): 6xxx and above: 24 numbers.}\textbf{==> Avoids redundant recounting by categorizing sets with minimum permutation > 2013 upfront, reducing per-set analysis effort}\\
Now let's add them up:\\
Set \{0,1,2,3\}: 5\\
Set \{1,2,3,4\}: 6\\
Set \{2,3,4,5\}: 24\\
Set \{3,4,5,6\}: 24\\
Set \{4,5,6,7\}: 24\\
Set \{5,6,7,8\}: 24\\
Set \{6,7,8,9\}: 24\\
Calculations: 5 + 6 = 11; 11 + 24*5 = 11 + 120 = \boxed{\textbf{131}}.\\
    \end{promptbox}
    \end{lrbox}
    \usebox{\DSRPrompt}
    \caption{The reasoner's response}
    \label{fig:evaluation-prompts}
\end{figure}




